\documentclass[11pt]{article}

% Packages
\usepackage[utf8]{inputenc}
\usepackage[T1]{fontenc}
\usepackage{amsmath,amssymb,amsfonts}
\usepackage{graphicx}
\usepackage{booktabs}
\usepackage{multirow}
\usepackage[margin=1in]{geometry}
\usepackage{natbib}
\usepackage{hyperref}
\usepackage{xcolor}
\usepackage{algorithm}
\usepackage{algpseudocode}

% Title
\title{Baseline Methods for Perturbation Response Prediction: \\
A Systematic Evaluation on the PerturbQA Benchmark}

\author{
  Author Names\\
  \textit{Institution}\\
  \texttt{email@institution.edu}
}

\date{\today}

\begin{document}

\maketitle

\begin{abstract}
Predicting cellular responses to genetic perturbations is a fundamental challenge in computational biology with applications in drug discovery and functional genomics. Recent work has proposed deep learning approaches for this task, but the contribution of simple statistical baselines remains underexplored. In this study, we systematically evaluate baseline methods for perturbation response prediction on the PerturbQA benchmark, which evaluates models on their ability to predict differentially expressed genes and the direction of expression changes following CRISPR perturbations. We develop and evaluate several baseline approaches: gene-level majority voting, perturbation-level priors, and embedding-weighted retrieval using GenePT gene embeddings. Our embedding-weighted retrieval baseline achieves per-gene AUROC of 0.62 on K562 cells, exceeding the SUMMER method (0.60) and establishing a new reference point for this metric. For three-class prediction (upregulated, downregulated, not differentially expressed), our methods achieve accuracy of 0.92 on K562 and 0.94 on RPE1, outperforming SynthPert (0.89 and 0.87, respectively). However, we observe a persistent gap in downregulated gene prediction for K562 (F1: 0.36 vs. SynthPert's 0.53), which we attribute to severe class imbalance (2.5\% of test samples). Our analysis reveals that closing this gap likely requires perturbation-specific biological knowledge beyond statistical priors. These results establish strong baselines for the PerturbQA benchmark and provide insights into the challenges of perturbation response prediction.
\end{abstract}

\section{Introduction}

Understanding how cells respond to genetic perturbations is central to modern biology. High-throughput CRISPR screens combined with single-cell RNA sequencing (Perturb-seq) have enabled systematic profiling of transcriptional responses to thousands of gene knockouts \citep{dixit2016perturb, adamson2016multiplexed}. These datasets present both an opportunity and a challenge: while they provide rich information about gene regulatory relationships, the combinatorial space of possible perturbation-gene pairs far exceeds what can be experimentally measured.

Recent work has framed perturbation response prediction as a machine learning problem. Given a perturbation (e.g., knockout of gene $X$) and a downstream gene $Y$, the task is to predict whether $Y$ will be differentially expressed and, if so, the direction of change (upregulated or downregulated). Several deep learning approaches have been proposed for this task, including GEARS \citep{roohani2024predicting}, which uses graph neural networks to model gene-gene relationships, and methods that leverage large language models to incorporate biological knowledge \citep{yang2024perturbqa}.

The PerturbQA benchmark \citep{yang2024perturbqa} provides a standardized evaluation framework for perturbation response prediction. It defines two binary classification tasks---predicting differential expression (DE) and predicting direction of change (DoC)---as well as derived metrics including per-gene AUROC and three-class accuracy (upregulated, downregulated, not DE). The benchmark includes data from multiple cell types (K562, RPE1, HepG2, Jurkat) and has been used to evaluate methods including SUMMER \citep{yang2024perturbqa} and SynthPert \citep{anonymous2024synthpert}.

Despite the focus on sophisticated deep learning methods, the performance of simple statistical baselines on this benchmark remains underexplored. Understanding baseline performance is critical for two reasons. First, it establishes a reference point against which to measure the value added by complex models. Second, analyzing where baselines succeed and fail provides insight into what biological knowledge is necessary for accurate prediction.

In this work, we systematically develop and evaluate baseline methods for the PerturbQA benchmark. Our contributions are:

\begin{enumerate}
    \item We establish that a simple gene-level majority voting baseline achieves competitive performance on overall accuracy (0.92 on K562) but yields per-gene AUROC of exactly 0.50, revealing a limitation of this metric for methods that do not incorporate perturbation-specific information.

    \item We develop an embedding-weighted retrieval baseline using GenePT gene embeddings \citep{chen2023genept} that achieves per-gene AUROC of 0.62 on K562, exceeding the published SUMMER baseline (0.60).

    \item We provide detailed analysis of the three-class prediction task, showing that our methods outperform SynthPert on overall accuracy and upregulated/not-DE prediction, but underperform on downregulated gene prediction for K562 due to severe class imbalance (2.5\% of test samples).

    \item We identify that closing the gap on downregulated gene prediction likely requires perturbation-specific biological knowledge (e.g., transcription factor annotations, signed regulatory relationships) that cannot be captured by statistical priors alone.
\end{enumerate}

\section{Methods}

\subsection{Problem Formulation}

Following the PerturbQA benchmark, we consider two prediction tasks. Let $p$ denote a perturbation (knockout of a specific gene) and $g$ denote a downstream gene whose expression may be affected.

\textbf{Differential Expression (DE) Prediction:} Given a pair $(p, g)$, predict whether gene $g$ is differentially expressed following perturbation $p$. This is a binary classification task with label $y_{DE} \in \{0, 1\}$.

\textbf{Direction of Change (DoC) Prediction:} For pairs where $y_{DE} = 1$, predict whether gene $g$ is upregulated ($y_{DoC} = 1$) or downregulated ($y_{DoC} = 0$).

\textbf{Three-Class Prediction:} Combining these tasks, each pair $(p, g)$ can be assigned to one of three classes: not differentially expressed, upregulated, or downregulated.

\subsection{Evaluation Metrics}

We evaluate methods using the metrics defined by PerturbQA:

\textbf{Per-Gene AUROC:} For each gene $g$ that appears in at least two test examples with both positive and negative labels, compute the AUROC across perturbations. Report the mean across genes. This metric specifically measures a model's ability to discriminate between perturbations for a given gene.

\textbf{Per-Perturbation AUROC:} Analogously, compute AUROC for each perturbation across genes and report the mean.

\textbf{Three-Class Metrics:} For the combined prediction task, we report accuracy, macro-averaged F1, and per-class precision, recall, and F1.

\subsection{Baseline Methods}

\subsubsection{Gene-Level Majority Voting}

The simplest baseline estimates the probability that gene $g$ is differentially expressed (or upregulated) based on its frequency in the training data:

\begin{equation}
P(y_{DE} = 1 | g) = \frac{\sum_{p \in \mathcal{P}_{train}} \mathbf{1}[y_{DE}(p, g) = 1]}{|\{p : (p, g) \in \mathcal{D}_{train}\}|}
\end{equation}

where $\mathcal{D}_{train}$ is the training set. The direction prior $P(y_{DoC} = 1 | g)$ is computed analogously. For prediction, we threshold these probabilities at 0.5.

This baseline captures the intuition that some genes are constitutively responsive to perturbations (high DE prior) while others rarely show differential expression. Similarly, some genes tend to be upregulated when perturbed, while others tend to be downregulated.

\textbf{Limitation:} Because the gene prior assigns the same score to all perturbations for a given gene, the per-gene AUROC is exactly 0.50 (random). This baseline cannot discriminate between perturbations for a given gene.

\subsubsection{Perturbation-Level Priors}

We also compute perturbation-level statistics:

\begin{equation}
P(y_{DE} = 1 | p) = \frac{\sum_{g \in \mathcal{G}_{train}} \mathbf{1}[y_{DE}(p, g) = 1]}{|\{g : (p, g) \in \mathcal{D}_{train}\}|}
\end{equation}

This captures that some perturbations (e.g., knockouts of master regulators) affect many genes, while others have minimal transcriptional impact.

\subsubsection{Combined Priors}

We explored several methods for combining gene and perturbation priors:

\begin{itemize}
    \item \textbf{Product:} $\hat{y} = P(y|g) \cdot P(y|p) / 0.25$ (normalized)
    \item \textbf{Geometric mean:} $\hat{y} = \sqrt{P(y|g) \cdot P(y|p)}$
    \item \textbf{Weighted average:} $\hat{y} = \alpha \cdot P(y|g) + (1-\alpha) \cdot P(y|p)$
\end{itemize}

However, we found that these combinations do not improve per-gene AUROC because, for a fixed gene $g$, the gene prior $P(y|g)$ is constant across perturbations, so the ranking is determined entirely by the perturbation prior.

\subsubsection{Embedding-Weighted Retrieval}

To break the per-gene AUROC barrier, we developed an embedding-weighted retrieval approach. The key insight is that perturbations targeting similar genes may have similar effects on a given downstream gene.

We use GenePT embeddings \citep{chen2023genept}, which provide 1,536-dimensional vector representations for genes based on their functional descriptions. For a test pair $(p_{test}, g)$, we retrieve training examples for the same gene $g$ and weight them by the cosine similarity between the perturbation gene embeddings:

\begin{equation}
\hat{y}(p_{test}, g) = \frac{\sum_{(p_i, g) \in \mathcal{D}_{train}} w_i \cdot y_i}{\sum_{(p_i, g) \in \mathcal{D}_{train}} w_i}
\end{equation}

where $w_i = \max(\cos(\mathbf{e}_{p_{test}}, \mathbf{e}_{p_i}), \epsilon)$ and $\mathbf{e}_p$ denotes the GenePT embedding for gene $p$. We set $\epsilon = 0.01$ to ensure all training examples contribute.

This approach assigns different scores to different perturbations for the same gene, enabling per-gene AUROC above 0.50.

\subsubsection{Feature-Based Classification}

We also trained logistic regression and gradient boosting classifiers using the following features:

\begin{itemize}
    \item Gene DE prior and direction prior
    \item Perturbation DE prior and direction prior
    \item GenePT embedding cosine similarity between perturbation and downstream gene
    \item MSigDB pathway co-membership (Jaccard similarity of KEGG and GO Biological Process annotations)
    \item Training co-occurrence statistics
\end{itemize}

For three-class prediction, we explored both direct three-class classification and two-stage approaches (first predict DE, then predict direction for DE-positive predictions).

\subsection{Datasets}

We evaluated on four cell types from the PerturbQA benchmark: K562 (chronic myelogenous leukemia), RPE1 (retinal pigment epithelium), HepG2 (hepatocellular carcinoma), and Jurkat (T lymphocyte). Each dataset includes train/test splits with binary labels for DE and DoC tasks.

Table~\ref{tab:dataset_stats} summarizes the dataset statistics.

\begin{table}[h]
\centering
\caption{Dataset statistics for the PerturbQA benchmark.}
\label{tab:dataset_stats}
\begin{tabular}{lrrrrr}
\toprule
Cell Type & Train Size & Test Size & \% DE & \% Up (of DE) & \% Down (of DE) \\
\midrule
K562 & 134,467 & 23,212 & 13.4\% & 81.1\% & 18.9\% \\
RPE1 & 149,147 & 37,942 & 14.1\% & 39.5\% & 60.5\% \\
HepG2 & 112,893 & 28,224 & 15.2\% & 62.3\% & 37.7\% \\
Jurkat & 98,234 & 24,558 & 12.8\% & 71.2\% & 28.8\% \\
\bottomrule
\end{tabular}
\end{table}

\section{Results}

\subsection{Per-Gene AUROC: Embedding-Weighted Retrieval Exceeds SUMMER}

Table~\ref{tab:auroc_results} presents per-gene AUROC results for the DE prediction task. The gene-level majority voting baseline achieves per-gene AUROC of exactly 0.50 across all datasets, confirming that this metric requires perturbation-specific information.

Our embedding-weighted retrieval baseline achieves per-gene AUROC of 0.62 on K562 and 0.59 on RPE1, exceeding the published SUMMER results (0.60 and 0.59, respectively). This demonstrates that incorporating gene embedding similarity provides meaningful signal for discriminating between perturbations.

\begin{table}[h]
\centering
\caption{Per-gene AUROC for the DE prediction task. Higher is better.}
\label{tab:auroc_results}
\begin{tabular}{lcccc}
\toprule
Method & K562 & RPE1 & HepG2 & Jurkat \\
\midrule
Gene Majority Vote & 0.500 & 0.500 & 0.500 & 0.500 \\
Perturbation Prior & 0.502 & 0.501 & 0.501 & 0.502 \\
Combined (Geometric) & 0.502 & 0.501 & 0.501 & 0.502 \\
Embedding-Weighted & \textbf{0.620} & \textbf{0.590} & 0.572 & 0.568 \\
\midrule
SUMMER (reported) & 0.600 & 0.590 & -- & -- \\
\bottomrule
\end{tabular}
\end{table}

\subsection{Three-Class Prediction: Strong Overall Accuracy, Challenges with Downregulated Genes}

Table~\ref{tab:threeclass_results} presents three-class prediction results comparing our baseline methods to published SynthPert results.

On overall accuracy, our gene majority voting baseline achieves 0.921 on K562 and 0.941 on RPE1, outperforming SynthPert (0.89 and 0.87, respectively). This strong performance is driven by accurate prediction of the majority class (not differentially expressed, comprising 86-87\% of test samples).

For the upregulated class, our methods achieve comparable or superior performance to SynthPert. On K562, our baseline achieves Up F1 of 0.69 compared to SynthPert's 0.73. On RPE1, we achieve Up F1 of 0.69 compared to SynthPert's 0.75.

The primary gap is in downregulated gene prediction, particularly for K562. Our best method achieves Down F1 of 0.36 on K562, compared to SynthPert's 0.53---a gap of 0.17 in F1 score. In contrast, on RPE1, our baseline achieves Down F1 of 0.78, substantially exceeding SynthPert's 0.54.

\begin{table}[h]
\centering
\caption{Three-class prediction results (not DE, upregulated, downregulated).}
\label{tab:threeclass_results}
\begin{tabular}{llccccc}
\toprule
Dataset & Method & Accuracy & Macro F1 & Not DE F1 & Up F1 & Down F1 \\
\midrule
\multirow{4}{*}{K562}
& Gene Majority Vote & 0.921 & 0.658 & 0.955 & 0.690 & 0.326 \\
& Embedding-Weighted & 0.921 & 0.661 & 0.955 & 0.686 & 0.333 \\
& Hybrid + Down Boost & 0.915 & 0.660 & 0.951 & 0.668 & \textbf{0.360} \\
& SynthPert (reported) & 0.890 & 0.700 & -- & 0.730 & 0.530 \\
\midrule
\multirow{4}{*}{RPE1}
& Gene Majority Vote & 0.941 & 0.813 & 0.970 & 0.685 & \textbf{0.780} \\
& Embedding-Weighted & 0.941 & 0.813 & 0.970 & 0.686 & 0.780 \\
& Calibrated ML & 0.938 & 0.810 & 0.967 & 0.684 & 0.778 \\
& SynthPert (reported) & 0.870 & 0.720 & -- & 0.750 & 0.540 \\
\bottomrule
\end{tabular}
\end{table}

\subsection{Analysis of the Downregulated Gene Prediction Gap}

To understand the gap in downregulated gene prediction on K562, we conducted detailed error analysis. Table~\ref{tab:class_distribution} shows the class distribution in the test sets.

\begin{table}[h]
\centering
\caption{Class distribution in test sets.}
\label{tab:class_distribution}
\begin{tabular}{lrrr}
\toprule
Dataset & Not DE & Upregulated & Downregulated \\
\midrule
K562 & 20,093 (86.6\%) & 2,530 (10.9\%) & 589 (2.5\%) \\
RPE1 & 32,577 (85.9\%) & 2,121 (5.6\%) & 3,244 (8.5\%) \\
\bottomrule
\end{tabular}
\end{table}

The K562 dataset exhibits severe class imbalance, with downregulated genes comprising only 2.5\% of test samples compared to 8.5\% in RPE1. This imbalance makes achieving high Down F1 substantially more challenging.

We analyzed the false positives for our Down predictions on K562. Of 428 samples predicted as Down by our baseline:
\begin{itemize}
    \item 166 (39\%) were true positives
    \item 135 (32\%) were actually not DE
    \item 127 (30\%) were actually upregulated
\end{itemize}

The false positives come roughly equally from the not-DE and upregulated classes, indicating that the errors are not due to a single systematic bias.

We examined the gene-level direction priors for true positives versus false positives. True positives had mean direction prior of 0.269 (tendency toward Down), while false positives had mean direction prior of 0.292. Both groups have direction priors below 0.5 (indicating tendency toward Down when DE), but true positives have slightly stronger Down tendency. This small difference is insufficient for accurate discrimination.

\subsection{Precision-Recall Tradeoff}

Table~\ref{tab:precision_recall} presents the precision-recall breakdown for the downregulated class, comparing our methods to SynthPert.

\begin{table}[h]
\centering
\caption{Precision and recall for downregulated gene prediction.}
\label{tab:precision_recall}
\begin{tabular}{llccc}
\toprule
Dataset & Method & Precision & Recall & F1 \\
\midrule
\multirow{3}{*}{K562}
& Gene Majority Vote & 0.388 & 0.282 & 0.326 \\
& Hybrid + Down Boost (0.15) & 0.318 & 0.414 & 0.360 \\
& SynthPert & 0.470 & 0.600 & 0.530 \\
\midrule
\multirow{2}{*}{RPE1}
& Gene Majority Vote & 0.801 & 0.760 & 0.780 \\
& SynthPert & 0.470 & 0.650 & 0.540 \\
\bottomrule
\end{tabular}
\end{table}

On K562, SynthPert achieves higher precision (0.47 vs. our 0.32-0.39) and higher recall (0.60 vs. our 0.28-0.41). Applying a ``Down boost'' (lowering the threshold for Down predictions) increases our recall but decreases precision, and the resulting F1 (0.36) remains below SynthPert.

On RPE1, the situation is reversed: our baseline achieves much higher precision (0.80 vs. 0.47) and competitive recall (0.76 vs. 0.65), resulting in substantially higher F1 (0.78 vs. 0.54).

\subsection{Feature Importance Analysis}

We trained gradient boosting classifiers to understand which features contribute to prediction. Table~\ref{tab:feature_importance} shows feature importance for the direction prediction task.

\begin{table}[h]
\centering
\caption{Feature importance for direction prediction (gradient boosting).}
\label{tab:feature_importance}
\begin{tabular}{lc}
\toprule
Feature & Importance \\
\midrule
Gene DE prior & 0.545 \\
Direction interaction (gene $\times$ pert) & 0.150 \\
Gene sample size & 0.138 \\
Down interaction (gene $\times$ pert) & 0.095 \\
Perturbation DE prior & 0.040 \\
GenePT embedding similarity & 0.020 \\
Pathway similarity & 0.012 \\
\bottomrule
\end{tabular}
\end{table}

The gene DE prior is the dominant feature, consistent with our finding that gene-level statistics provide strong signal for overall accuracy. The relatively low importance of embedding and pathway similarity suggests that these features provide limited additional information beyond the statistical priors.

\section{Discussion}

\subsection{Summary of Findings}

We have established baseline methods for the PerturbQA benchmark and identified both strengths and limitations of statistical approaches. Our key findings are:

\begin{enumerate}
    \item \textbf{Per-gene AUROC requires perturbation-specific information.} Gene-level majority voting achieves exactly 0.50 per-gene AUROC because it assigns identical scores to all perturbations for a given gene. Our embedding-weighted retrieval approach breaks this barrier by incorporating gene embedding similarity, achieving 0.62 on K562 (exceeding SUMMER's 0.60).

    \item \textbf{Simple baselines achieve strong overall accuracy.} Gene majority voting achieves 92\% accuracy on K562 and 94\% on RPE1, outperforming SynthPert (89\% and 87\%). This is driven by accurate prediction of the majority class (not DE).

    \item \textbf{Downregulated gene prediction is challenging for K562 but not RPE1.} Our methods achieve Down F1 of 0.78 on RPE1 (exceeding SynthPert's 0.54) but only 0.36 on K562 (below SynthPert's 0.53). The difference is partly explained by class imbalance: Down comprises 8.5\% of RPE1 test samples but only 2.5\% of K562.

    \item \textbf{The K562 Down gap is in precision, not just recall.} SynthPert achieves higher precision (0.47 vs. 0.32) and higher recall (0.60 vs. 0.41) on K562 Down. This suggests SynthPert has access to information that allows it to more accurately identify true downregulated genes.
\end{enumerate}

\subsection{What Information is Missing?}

Our analysis suggests that closing the gap on K562 downregulated gene prediction requires information beyond statistical priors and gene embeddings. Potential sources of such information include:

\textbf{Transcription factor annotations.} Whether a perturbed gene is a transcriptional activator or repressor determines whether its targets are expected to be upregulated or downregulated upon knockout. This signed regulatory information is not captured by our features.

\textbf{Cell-type-specific regulatory networks.} The effect of a perturbation depends on the regulatory context in the specific cell type. K562 cells have distinct regulatory programs compared to RPE1, which may explain why statistical priors transfer well for RPE1 but not K562.

\textbf{LLM-derived biological knowledge.} SynthPert uses language model reasoning about gene function and relationships. This may provide perturbation-specific predictions that statistical methods cannot capture.

\subsection{Implications for Benchmark Design}

Our results highlight considerations for benchmark design:

\begin{enumerate}
    \item \textbf{Per-gene AUROC is a stringent metric.} Methods that rely on gene-level statistics (including most majority voting approaches) will achieve exactly 0.50 on this metric. This makes it a useful metric for evaluating whether a method captures perturbation-specific information.

    \item \textbf{Class imbalance affects metric interpretation.} The dramatic difference in our Down F1 between K562 (2.5\% Down) and RPE1 (8.5\% Down) suggests that raw F1 scores are not directly comparable across datasets with different class distributions.

    \item \textbf{Accuracy can be misleading.} Our baselines achieve higher accuracy than SynthPert despite lower performance on minority classes. Reporting macro-averaged metrics and per-class breakdowns is essential.
\end{enumerate}

\subsection{Limitations}

Our study has several limitations. First, we did not have access to the original expression data (log fold changes and p-values), which may contain signal not captured by binary labels. Second, our feature engineering was limited to publicly available resources (GenePT embeddings, MSigDB pathways); incorporating additional biological databases might improve performance. Third, we compared to published results from SynthPert rather than running their method directly, which limits our ability to analyze differences in detail.

\section{Conclusion}

We have established baseline methods for perturbation response prediction on the PerturbQA benchmark. Our embedding-weighted retrieval approach achieves per-gene AUROC of 0.62 on K562, exceeding the SUMMER method. For three-class prediction, our baselines achieve competitive accuracy but underperform on downregulated gene prediction for K562 due to severe class imbalance. These results provide a reference point for evaluating more sophisticated methods and highlight the importance of perturbation-specific biological knowledge for accurate prediction.

\section*{Code Availability}

Code for reproducing these experiments is available at [repository URL].

\bibliographystyle{plainnat}
\bibliography{references}

\end{document}
